\begin{enumerate}

\item Las áreas son $A_1 = 2400$ mm$^2$ y $A_2 = 2000$ mm$^2$, con centros a $10$ y $70$ mm de la cara superior.

\item El centro de gravedad es la media ponderada: $y_G = (A_1 \cdot 10 + A_2 \cdot 70)/(A_1 + A_2)$.

\item Cada rectángulo aporta $b\,h^3/12$ propio más $A_i d_i^2$ por la distancia de su centro al centro de gravedad.

\end{enumerate}

\noindent\textbf{Operaciones y resultados intermedios.}

\begin{enumerate}

\item Posición del centro de gravedad.

\[\fitmath{y_{1}=\frac{2400\cdot 10+2000\cdot 70}{2400+2000}\approx 37.27273\,\mathrm{mm}}\]

\item Momento de inercia del conjunto.

\[\fitmath{y_{2}=\frac{120\cdot {20}^{3}}{12}+2400\cdot {\left(\frac{2400\cdot 10+2000\cdot 70}{4400}-10\right)}^{2}+\frac{20\cdot {100}^{3}}{12}+2000\cdot {\left(70-\frac{2400\cdot 10+2000\cdot 70}{4400}\right)}^{2}\approx 5673939\,\mathrm{mm^{4}}}\]

\end{enumerate}
